\documentclass{article}
\usepackage{hyperref}
\usepackage{Style}

\nocite{*} % Comentar si quiero citar
%\addbibresource{bibliografia.bib} % Quitar el comentado si quiero usar bibliografia

\begin{document}

\begin{minipage}{2.5cm}
    \includegraphics[width=2cm]{imagen_puc.jpg}
\end{minipage}
\begin{minipage}{12cm}
    {\sc Pontificia Universidad Católica de Chile\\
    Facultad de Matemáticas\\
    Departamento de Matemática\\
    Profesor: Manuel Cabezas -- Ayudante: Benjamín Mateluna}
\end{minipage}
\vspace{2ex}

{\centerline{\bf Introducción a la Geometría - MAT1304}
\centerline{\bf Ayudantía 17 - Repaso Midterm}}
\centerline{\bf 07 de mayo de 2026}

\vspace{1cm}
\noindent\textbf{Problema 1.} Sean $A,B$ y $C$ puntos tales que $\angle ABC=\alpha$, demuestre que
existe un punto $D$ tal que $\angle ABD=2\alpha$.

\vspace{5mm}
\noindent\textbf{Problema 2.} Sea $\triangle ABC$ y $x\in\R$ positivo. Supongamos que 
$a=2x,b=x+1,c=x+2$ y que $\alpha=\frac{\pi}{3}$. Pruebe que $\beta$ es constructible con regla y 
compás.

\vspace{5mm}
\noindent\textbf{Problema 3.} En $\triangle ABC$, se cumple que
\begin{equation*}
    \frac{a}{cos(\alpha)}=\frac{b}{cos(\beta)}=\frac{c}{cos(\gamma)}
\end{equation*}
Demuestre que el triángulo es equilátero.

\vspace{5mm}
\noindent\textbf{Problema 4.} Encuentre polinomios $q(x)$ y $p(x)$ de grado $2$ tales que
\begin{equation*}
    q(x)p(x)=x^{4}+3x^{3}+4x^{2}+3x+1
\end{equation*}

\vspace{5mm}
\noindent\textbf{Problema 5.} Muestre que $cos(\frac{2\pi}{9})$ es irracional.

\vspace{1cm}
\noindent\textbf{\underline{Ejercicios Propuestos:}}

\vspace{5mm}
\noindent\textbf{Extra 1.} Demuestre la siguiente identidad trigonométrica
\begin{equation*}
    tan\left(\frac{\alpha}{2}\right)=csc(\alpha)-cot(\alpha)
\end{equation*}

\vspace{5mm}
\noindent\textbf{Extra 2.} Calcular el valor de $sen(\frac{\pi}{5})$ y uselo para calcular 
$sen(\frac{\pi}{20})$.

%\printbibliography % Quitar el comentado si quiero usar bibliografia

\end{document}